% Standalone section fragment; theorem environments supplied by parent manuscript.
\section{A cancellation-resistant tensor detector for icosahedral fields}
\label{sec:artin}

This section supplies a separate transfer interface. It does not improve the
smallest detector already in the incoming corpus: the $q=4$ case of the prior
$\mathrm{PSL}_2(q)$ report already has a degree-$16$ detector. The additional calculation
here is the exact first tensor power of a specified three-dimensional
representation for which proper-subgroup information alone certifies
holomorphy and detects every Dedekind zero, together with an optimal
multiplicity certificate. No literature-priority claim is made for these
finite-character identities.

\subsection{The general tensor construction}
Let $E/K$ be finite Galois with group $G$. At a point $s_0\ne1$, put
\[
 n_\psi=\operatorname{ord}_{s_0}L(s,\psi,E/K),\qquad
 \Theta=\sum_{\psi\in\operatorname{Irr}(G)}n_\psi\psi,\qquad
 d=\Theta(1)=\operatorname{ord}_{s_0}\zeta_E.
\]
The orders exist by Brauer's meromorphic continuation theorem; poles have
negative order. Classical Artin formalism gives
\[
 \langle\operatorname{Res}_C\Theta,\lambda\rangle
 =\operatorname{ord}_{s_0}L(s,\lambda,E/E^C)
 \quad(C\le G,\ \lambda\in\operatorname{Irr}(C)).
\]
When $C$ is cyclic the right side is nonnegative: its factors are finite-order
Hecke functions, with the sole possible pole at $1$. Hence
$\operatorname{Res}_C\Theta$ is a genuine character, $d\ge0$, and
$|\Theta(g)|\le d$ for every $g\in G$. Moreover $d=0$ forces $\Theta=0$.
These are classical Heilbronn-character facts, not consequences of the
proposed RH proof; see Foote--Ginsberg--Murty and Gun--Hazra--Sahu \cite{heilbronn,gun}.

Assume $G$ is nontrivial; the trivial case is immediate.
Let $\rho$ be a faithful complex representation of dimension $r$, and set
$\tau=1\oplus\rho$, $D=r+1$. Its character has
$|\chi_\tau(g)|<D$ for $g\ne1$: equality in the triangle inequality would
force every eigenvalue, including the distinguished eigenvalue $1$, to be
$1$, contradicting faithfulness. Define
\[
 F_m(g)=|\chi_\tau(g)|^{2m},\quad
 q=\max_{g\ne1}\frac{|\chi_\tau(g)|^2}{D^2}<1,
 \quad
 \kappa_m=\frac{F_m(1)-\sum_{g\ne1}F_m(g)}{|G|}.
\]
Here $F_m$ is the genuine character of
$(\tau\otimes\overline\tau)^{\otimes m}$ and is nonnegative pointwise.
If $(|G|-1)q^m<1$, then $\kappa_m>0$ and
\begin{equation}\label{eq:generic-tensor-domination}
 \kappa_m d\le\operatorname{ord}_{s_0}L(s,F_m,E/K)
 \le\langle F_m,1\rangle d.
\end{equation}
Indeed the inner product with $\Theta$ is real, its identity contribution is
$F_m(1)d/|G|$, and all other terms have absolute value at most
$F_m(g)d/|G|$. The upper bound follows by summing the same pointwise bound.
Consequently $L(s,F_m)$ is holomorphic away from $1$ and has precisely the
same zero support as $\zeta_E$ away from $1$. Its pole order at $1$ is
$\langle F_m,1\rangle$. This establishes holomorphy for an aggregate tensor
representation, not for the individual constituents of $\rho$.

For an explicit group-only sufficient exponent, if $e$ is the exponent of
$G$ then
\[
 q\le1-\frac{4\sin^2(\pi/e)}{D^2}.
\]
To prove this, diagonalize $\rho(g)$ unitarily. At least one eigenvalue
$\lambda\ne1$ is an $e$th root of unity, and
$D^2-|\chi_\tau(g)|^2$ is the sum of the squared pairwise distances among
all $D$ eigenvalues. The pair $1,\lambda$ contributes at least
$4\sin^2(\pi/e)$. This bound is deliberately crude. The character-table
certificate below is much sharper. Tensor generation itself is classical;
compare the Burnside--Brauer discussion in Steinberg \cite{steinberg}.

\subsection{An exact fourth-tensor certificate for \texorpdfstring{$A_5$}{A5}}
Write $\chi,\chi',\eta,\nu$ for the irreducible characters of degrees
$3,3,4,5$. Let $\phi=(1+\sqrt5)/2$, $\phi'=(1-\sqrt5)/2$. The table is
\[
\begin{array}{c|rrrrr}
 &1&2A&3A&5A&5B\\\hline
\text{class size}&1&15&20&12&12\\
1&1&1&1&1&1\\
\chi&3&-1&0&\phi&\phi'\\
\chi'&3&-1&0&\phi'&\phi\\
\eta&4&0&1&-1&-1\\
\nu&5&1&-1&0&0
\end{array}
\]
Let $D_5$ denote a dihedral subgroup of order $10$, with nontrivial sign
character $\epsilon$ trivial on its rotations. Choose $C_5\le D_5$ and a
nontrivial character $\lambda$ corresponding to the eigenvalues of $\chi$.
Let $\omega$ be either nontrivial character of $A_4/V_4$. Frobenius
reciprocity, or the explicit induction calculation supplied with the
archive, gives
\[
 A=\operatorname{Ind}_{D_5}^{A_5}\epsilon=\chi+\chi',\qquad
 B=\operatorname{Ind}_{C_5}^{A_5}\lambda=\chi+\eta+\nu,\qquad
 C=\operatorname{Ind}_{A_4}^{A_5}\omega=\nu.
\]
The tensor powers decompose as
\[
 \chi^2=1+\chi+\nu,\qquad
 \chi^3=1+3\chi+\chi'+\eta+2\nu,\qquad
 F=\chi^4=3+6\chi+3\chi'+4\eta+7\nu.
\]
The decisive positive identity is
\begin{equation}\label{eq:a5-fourth-certificate}
 \boxed{\ 4F=\operatorname{Reg}_{A_5}+11\cdot1+9A+12B+11C.\ }
\end{equation}
All coefficients on the right are nonnegative. For every $A_5$ extension
$E/\mathbb Q$, Artin formalism therefore gives the exact meromorphic identity
\begin{equation}\label{eq:a5-artin-product}
 L(s,F)^4=\zeta_E(s)\zeta_{\mathbb Q}(s)^{11}
 L(s,\epsilon,E/E^{D_5})^9
 L(s,\lambda,E/E^{C_5})^{12}
 L(s,\omega,E/E^{A_4})^{11}.
\end{equation}
The three displayed nontrivial one-dimensional factors are entire Hecke
functions and are nonzero at $s=1$. The right-hand side therefore has
pole order $12$ there. Since $L(s,F)$ is already meromorphic, its fourth power being
holomorphic away from $1$ implies that it too is holomorphic there.
It has pole order $3$ at $1$, the multiplicity of the trivial character.

At any $s_0\ne1$, set $e_F=\operatorname{ord}_{s_0}L(s,F)$ and
$d=\operatorname{ord}_{s_0}\zeta_E$. Equation
\eqref{eq:a5-artin-product} proves $e_F\ge d/4$.
Since $F(g)=\chi(g)^4\ge0$ and $|\Theta(g)|\le d$,
\[
 e_F=\frac1{60}\sum_gF(g)\overline{\Theta(g)}
 \le\frac d{60}\sum_gF(g)=3d.
\]
Thus
\begin{equation}\label{eq:a5-multiplicities}
 \boxed{\quad d/4\le e_F\le3d.\quad}
\end{equation}
This proves identical zero support, even when individual Artin factors
might have poles. It is stronger than merely observing that every
irreducible occurs in $F$; positive constituent multiplicities by
 themselves do not prevent cancellation of meromorphic orders.

\subsection{Exact scope of the optimality statement}
Consider the integral virtual character
\[
 \Theta_*=-\chi+\chi'+\eta.
\]
Its degree is $4$, and its restrictions to representatives of all three
maximal-subgroup types are genuine:
\[
\begin{array}{c|l}
 A_4&1+\text{the three-dimensional irreducible}\\
 D_5&2\cdot\text{one two-dimensional irreducible}\\
 S_3&1+\operatorname{sgn}+\text{the two-dimensional irreducible}.
\end{array}
\]
Each of these maximal subgroups is monomial: its irreducible characters
are induced from linear characters, so Artin--Hecke theory really does
impose these nonnegative restriction orders at an arithmetic point.
Every proper subgroup is contained in one of these maximal subgroups.
The included permutation enumeration verifies the $59$ subgroups of $A_5$
and its maximal subgroup counts: ten of order $6$, six of order $10$, and
five of order $12$. Hence $\Theta_*$ satisfies all necessary
proper-subgroup holomorphy constraints, including every cyclic restriction.
However
\[
 \langle\chi^j,\Theta_*\rangle=-1\ (j=1,2,3),\qquad
 \langle\chi^4,\Theta_*\rangle=1.
\]
Therefore the first three tensor powers cannot be certified by those
constraints, and no inequality $d\le k e_F$ valid throughout that cone can
have $k<4$. The identity \eqref{eq:a5-fourth-certificate} attains $k=4$.
This proves optimality for the algebraic cone of necessary character
conditions. It does not assert that an actual zero of any number field
realizes $\Theta_*$, nor that the first three tensor functions have poles.

For a complete small cone computation, write
$\Theta=n_0\cdot1+x\chi+y\chi'+z\eta+w\nu$ and suppose $n_0=0$.
The genuine restrictions to $A_4,D_5,S_3$ are equivalent to
\[
 z\ge0,\quad w\ge0,\quad x+y\ge0,\quad x+z+w\ge0,\quad y+z+w\ge0.
\]
Its six extreme rays in coordinates $(x,y,z,w)$ are
\[
 (1,0,0,0),\ (0,1,0,0),\ (-1,1,1,0),\ (1,-1,1,0),\
 (-1,1,0,1),\ (1,-1,0,1).
\]
To see completeness, put $s=x+y\ge0$ and
$t=(x+z+w)/(s+2z+2w)\in[0,1]$ for a nonzero vector. The vector is the
convex combination, with coefficients $t,1-t$, of
$s(1,0,0,0)+z(1,-1,1,0)+w(1,-1,0,1)$ and
$s(0,1,0,0)+z(-1,1,1,0)+w(-1,1,0,1)$.
The ratios $e_F/d$ on these rays are respectively
$2,1,1/4,7/4,4/5,2$. Thus the exact bounds in this cone are
$d/4\le e_F\le2d$. Under ordinary RH, they apply throughout
$1/2<\Re s_0<1$.

The prior degree-$16$ detector is $F_0=1+\chi+\chi'+\eta+\nu$,
with class values $(16,0,1,1,1)$ \cite{pslprior}. It can also be written
$1+\operatorname{Ind}_{V_4}^{A_5}\lambda_0$ for a nontrivial linear
character $\lambda_0$ of $V_4$. This is a more economical detector when
degree is the objective; it was already present in the inspected corpus.

For comparison, the crude pointwise dominance test fails for $\chi^4$:
its identity value is $81$, while its nonidentity sum is
$15+12(\phi^4+\phi'^4)=99$. It first succeeds among the even powers at
$\chi^6$, where $729>15+12(\phi^6+\phi'^6)=231$.
The proper-subgroup certificate therefore removes one tensor-square
amplification step; it does not improve the prior degree-$16$ detector.

\subsection{The analytic estimate still needed}
For a ramified prime, let $I$ be inertia and let $g$ be a Frobenius lift in
the decomposition group. Averaging the inertia projector gives
\[
 \operatorname{tr}\bigl(g^k\mid(\rho^{\otimes4})^I\bigr)
 =\frac1{|I|}\sum_{h\in I}\chi(g^kh)^4\ge0.
\]
Consequently the coefficients $b(p^k)$ of $D_F=-L'/L(s,F)$ are real,
nonnegative, and at most $81\log p$, including ramified primes. Let
\[
 D_F(s)=\frac3{s-1}+a_0+O(s-1)
\]
and define
\[
 T_F(x)=\sum_{n\le x}b(n)(x/n-1)-3x\log x-(a_0-3)x.
\]
For $\Re s>1$, direct integration gives
\[
 \int_1^\infty T_F(x)x^{-s-1}\,dx
 =\frac{D_F(s)}{s(s-1)}-\frac3{(s-1)^2}-\frac{a_0-3}{s-1}.
\]
The singularity at $1$ cancels. A zero $\rho$ of $L(s,F)$ contributes
residue $e_F/[\rho(1-\rho)]$, strictly positive for a real zero
$0<\rho<1$. Thus the same one-sided Mellin--Landau principle used in the
source would imply Dedekind RH for $E$ from
\[
 T_F(x)\le C_{E,\varepsilon}x^{1/2+\varepsilon}
 \quad\text{for every }\varepsilon>0\text{ and all sufficiently large }x.
\]
No such arithmetic bound is established here. The new burden is to build
short and sufficiently spanning integral sections whose long-vector
estimate encodes this tensor trace, with compatible twist and
ramification bounds. Character identities and holomorphy alone do not
supply that construction. In particular, this section does not solve
individual Artin holomorphy or establish a new case of Dedekind RH.
